% Generated draft: requires mathematical and source review.
\documentclass[11pt]{article}
\usepackage[margin=25mm]{geometry}
\usepackage{fontspec}
\setmainfont{Latin Modern Roman}
\usepackage{amsmath,amssymb,mathtools,unicode-math}
\setmathfont{Latin Modern Math}
\usepackage{enumitem,needspace,xurl,hyperref}
\setlength{\parindent}{0pt}
\setlength{\parskip}{5pt}
\setlength{\emergencystretch}{3em}
\newcommand{\N}{\mathbb{N}}
\newcommand{\Z}{\mathbb{Z}}
\newcommand{\Q}{\mathbb{Q}}
\newcommand{\R}{\mathbb{R}}
\newcommand{\C}{\mathbb{C}}
\newcommand{\F}{\mathbb{F}}
\newcommand{\A}{\mathbb{A}}
\newcommand{\PP}{\mathbb{P}}
\newcommand{\E}{\mathbb{E}}
\newcommand{\eps}{\varepsilon}
\newcommand{\dd}{\,\mathrm{d}}
\begin{document}
\section*{UnsolvedMath ID: 30005369}
\textbf{Source problem code:} catalogue entry
\subsection{Definitions and mathematical statement}
Under the conventional smooth interpretation, a knot is a smooth oriented embedding of \(S^1\) in \(S^3\), considered up to oriented ambient isotopy. A knot is \emph{positive} if it admits an oriented diagram in which every crossing is positive, using the usual right-handed crossing convention; a crossing-free diagram is allowed. Two knots \(K_0,K_1\subset S^3\) are \emph{smoothly concordant} if there is a smoothly and properly embedded oriented annulus \(A\subset S^3\times[0,1]\) whose oriented boundary is
\[
\partial A=(-K_0)\times\{0\}\;\sqcup\;K_1\times\{1\}.
\]
They are \emph{oriented ambient isotopic} if a smooth ambient isotopy \(H_t:S^3\to S^3\), with \(H_0=\operatorname{id}_{S^3}\), carries \(K_0\) to \(K_1\) with its orientation.

\emph{Question under this interpretation.} Is it true that, for every pair of positive knots \(K_0,K_1\subset S^3\),
\[
K_0\text{ and }K_1\text{ are smoothly concordant}
\quad\Longrightarrow\quad
K_0\text{ and }K_1\text{ are oriented ambient isotopic}?
\]
The supplied source record does not explicitly specify the ambient manifold, orientation conventions, or whether concordance and isotopy are taken in the smooth or topological category. The displayed statement therefore makes a conventional smooth interpretation explicit, pending source confirmation. A version using locally flat topological concordance is a distinct formulation and is not asserted to be equivalent.
\subsection{Short English statement}
Must any two positive knots that are concordant be isotopic? The displayed formulation uses smooth oriented knots in \(S^3\); these conventions require confirmation from the source. The supplied record labels the problem open, but that status has not been independently verified.
\subsection{Sources}
\begin{itemize}
\item Local UnsolvedMath registry, permanent ID \(30005369\), source problem code \texttt{catalogue entry}, title \emph{Concordance Versus Isotopy for Positive Knots}. The supplied record belongs to the Oberwolfach Reports Open Problems collection and assigns the year 2023; no individual proposer is identified.
\item Supplied primary-source citation: \emph{Low-dimensional Topology} (2023), Oberwolfach Reports, DOI 10.4171/owr/2023/3, pp.\ 151--213. \url{https://doi.org/10.4171/owr/2023/3}. The registry describes these reports as documenting mathematical workshops and their open-problem discussions.
\item Supplied literature source, labelled ``Baader--Dehornoy--Liechti, positive knots'': \url{https://arxiv.org/abs/1503.01946}.
\item Literature assessment dated \texttt{2026-08-22}: the supplied curation labels the problem open and reports that primary discussions still treat concordance-implies-isotopy for positive knots as unresolved, while existing results provide finiteness or special-family theorems rather than a general resolution. This is supplied dated curation, not an independently verified literature assessment. Its listed sources are the Oberwolfach report and the Baader--Dehornoy--Liechti link above; the record requests independent verification of time-sensitive status and cited claims.
\end{itemize}
\end{document}
